\begin{figure}[htbp]
\centering
\begin{tikzpicture}[x=.82cm,y=.82cm,font=\small]
  \fill[accent!22] (0,0) rectangle (3,3);
  \fill[white] (1,1) rectangle (2,2);
  \foreach \x in {0,2} \foreach \y in {0,2}
    {\fill[accent!52] (\x,\y) rectangle ++(1,1);}
  \draw[ink] (0,0) rectangle (3,3);
  \draw[ink] (1,1) rectangle (2,2);
  \foreach \x in {1,2} {\draw[ink!60] (\x,0)--(\x,3);}
  \foreach \y in {1,2} {\draw[ink!60] (0,\y)--(3,\y);}
  \node at (1.5,.5) {1}; \node at (2.5,1.5) {2};
  \node at (1.5,2.5) {3}; \node at (.5,1.5) {4};
  \node[anchor=north] at (1.5,-.15) {Four closed strips};
  \begin{scope}[shift={(6.5,1.5)}]
    \coordinate (a) at (0,-1.15); \coordinate (b) at (1.15,0);
    \coordinate (c) at (0,1.15); \coordinate (d) at (-1.15,0);
    \draw[accent,thick] (a)--(b)--(c)--(d)--cycle;
    \foreach \pt in {a,b,c,d} {\fill[ink] (\pt) circle (2.5pt);}
    \node[below] at (a) {1}; \node[right] at (b) {2};
    \node[above] at (c) {3}; \node[left] at (d) {4};
    \node[anchor=north] at (0,-1.65) {Intersection nerve};
  \end{scope}
\end{tikzpicture}
\caption{Each darker square is an intersection of two strips and gives one edge of the nerve. There are no triple intersections. Both spaces have one independent loop.}
\label{fig:frame}
\end{figure}
